\subsection{Phân Tích Phần Dư}


Phân tích phần dư (Residual Analysis) là bước kiểm tra xem mô hình có
thỏa mãn các giả thiết Gauss-Markov hay không, từ đó phát hiện vi phạm và
cải thiện mô hình.

\textbf{Cài đặt Python - F8: \texttt{residual\_plots}} tạo ra 4 biểu đồ
chẩn đoán chuẩn:

\begin{enumerate}
    \item \textbf{Residuals vs Fitted:} Kiểm tra tính tuyến tính (GM1) và
          đồng phương sai (GM4). Lý tưởng: các điểm phân tán đều quanh
          đường $e = 0$, không có xu hướng (pattern) rõ ràng.

    \item \textbf{Normal Q-Q Plot:} So sánh quantile của phần dư với
          quantile của phân phối chuẩn lý thuyết (GM5). Lý tưởng: các điểm
          nằm trên đường chéo.

    \item \textbf{Scale-Location} ($\sqrt{|e_{\mathrm{std}}|}$ vs Fitted):
          Kiểm tra homoscedasticity (GM4). Lý tưởng: đường LOWESS nằm ngang.

    \item \textbf{Cook's Distance:} Phát hiện các quan sát có ảnh hưởng
          lớn (influential points) dựa trên leverage $h_{ii}$ và phần dư:
          \begin{equation}
              D_i = \frac{e_i^2\, h_{ii}}{(p+1)\hat{\sigma}^2(1 - h_{ii})^2}.
              \label{eq:cooks_d}
          \end{equation}
          Quy tắc: quan sát $i$ được coi là influential nếu $D_i > 4/n$.
\end{enumerate}

\begin{tcolorbox}[mycodebox, title={\textbf{Lưu ý triển khai}}]
Hàm \texttt{residual\_plots} gọi \texttt{hat\_matrix} (F2) để tính
leverage $h_{ii} = H_{ii}$ khi \texttt{X} được truyền vào, đảm bảo
Cook's Distance được tính chính xác với đúng số bậc tự do $n - p - 1$.
Khi \texttt{X=None}, hàm fallback sang $|e_i|$ (đủ cho hình 1, 2, 3 nhưng
hình 4 sẽ kém chính xác).
\end{tcolorbox}

\begin{figure}[h]
    \centering
    \includegraphics[width=\textwidth]{assets/part1/residual_plots.png}
    \caption{Bốn biểu đồ chẩn đoán phần dư. Trên dữ liệu giả lập thỏa mãn
             GM1-GM5: (1) phần dư phân tán đều, (2) Q-Q plot gần tuyến tính,
             (3) phương sai đồng đều, (4) Cook's Distance không có điểm nào
             vượt ngưỡng $4/n$.}
    \label{fig:residual_plots}
\end{figure}


